% =====================================================================
% DRAFT — SQ1: does the sampling segment of the AIF matter?
% Methods -> sections/02_methods.tex, Results -> sections/03_results.tex,
% Discussion -> sections/04_discussion.tex. References: methods_aif_refs.bib.
%
% Figures/tables to copy into the report:
%   main text   Table~\ref{tab:sq1_main}     (this file, numbers from results/combined_v1/sq1_full_vs_top/sq1_tests.csv)
%               Table~\ref{tab:sq1_context}  (this file, from sq1_report.md "Mask context" + plane_difference/plane_tests.csv)
%               Fig.~\ref{fig:sq1_forest}    results/combined_v1/sq1_full_vs_top/forest.png -> figures/sq1_forest.png
%   appendix D  full table                   results/combined_v1/sq1_full_vs_top/table_sq1.tex (all 15 features)
%               Bland--Altman                ba_grown.png, ba_model.png -> figures/sq1_ba_grown.png, figures/sq1_ba_model.png
%               Q--Q plots                   qq_grown.png, qq_model.png
%
% All numbers: QC with tail threshold 0.30 (current pipeline). If the 0.15 threshold is adopted, re-run
% aif_qc.py -> aif_taxonomy.py -> aif_feature_set.py -> sq1_full_vs_top.py and update the numbers.
% =====================================================================

% ---------------------------------------------------------------------
% METHODS
% ---------------------------------------------------------------------
\subsection{Sampling segment (SQ1)}
\label{sec:sq1_methods}

perfDSA samples the AIF from the supraclinoid ICA, chosen because it is easy to identify and less affected by
overprojected contrast and by contrast stasis than the lower ICA~\cite{su2025perfdsa}. In perfusion imaging in
general, the AIF site matters because bolus delay and dispersion between the sampling site and the tissue bias
CBF and MTT~\cite{calamante2000,calamante2013}, and in CT perfusion moving the AIF changed the CBF and MTT maps
markedly~\cite{ferreira2010}. SQ1 therefore asks whether the AIF, and the perfusion parameters derived from it,
change when the AIF is sampled from a longer stretch of the ICA instead of the supraclinoid segment alone.

Two full-ICA AIFs were compared with the supraclinoid \texttt{ica\_top} AIF of the same run:
\texttt{ica\_full\_model}, segmented by a network trained on the full ICA, and \texttt{ica\_full\_grown}, grown from
the \texttt{ica\_top} mask along the opacified ICA into the petrous and cervical segments. A pair was formed only
when both curves of the run were trusted (Section~\ref{sec:aif_qc}), so that differences reflect the segment and
not a failed extraction. Because the projection determines which parts of the ICA overlap with other vessels,
all analyses were done separately for plane A (frontal) and plane B (lateral).

For each pair the difference full $-$ top was computed for the timing, amplitude and shape features of
Section~\ref{sec:features} and for the median CBV, CBF, MTT and Tmax over the perfused pixels: as a plain
difference for times and dimensionless shape features, and as a log ratio for amplitudes and perfusion
parameters, whose error scales with their size (reported as \%). Normality of the differences was assessed with
the Shapiro--Wilk test and Q--Q plots. A paired $t$-test was used when $p_\mathrm{SW} \geq 0.05$, otherwise the
Wilcoxon signed-rank test, which assumes only a symmetric distribution of the differences; both $p$-values were
retained to check that the conclusions do not depend on this choice. $P$-values were Holm-corrected over the 15
features within each comparison and plane. Agreement was quantified with Bland--Altman analysis (bias and 95\,\%
limits of agreement, LoA). Temporal agreement of the two curves was summarised with the Pearson correlation.

To test whether the effect of the segment depends on the view, the difference between planes ($d_B - d_A$) was
tested per injection with the Wilcoxon signed-rank test (Holm-corrected). Three candidate explanations were
related to the full $-$ top difference with Spearman correlation: the number of averaged pixels (log ratio of the
mask sizes), vessel overlap (difference between the masks in the fraction of vessel pixels in a 15-pixel ring
around the mask, and in the late-opacifying fraction) and foreshortening of the supraclinoid segment (size and
spatial spread of the \texttt{ica\_top} mask). Finally, the AIF of each variant was normalised to the whole-image
contrast signal of its run (Section~\ref{sec:plane_methods}) to determine which of the two AIFs changes between
the views.

% ---------------------------------------------------------------------
% RESULTS
% ---------------------------------------------------------------------
\subsection{Sampling segment (SQ1)}
\label{sec:sq1_results}

\textbf{Pairs.} Both curves were trusted in 44 (A) and 51 (B) runs for \texttt{ica\_full\_model} and in 46 (A)
and 53 (B) runs for \texttt{ica\_full\_grown}. Shapiro--Wilk rejected normality for most
\texttt{ica\_full\_model} differences (normal in 0/15 features in A, 2/15 in B): the two masks are nearly the
same, so the differences are concentrated around zero with a few larger values. For \texttt{ica\_full\_grown},
8/15 (A) and 10/15 (B) were normal. The paired $t$-test and the Wilcoxon test disagreed on significance in
7 of 60 tests, none of them among the results reported below as large.

\textbf{Model-segmented full ICA.} The \texttt{ica\_full\_model} mask covered almost the same pixels as
\texttt{ica\_top} (median Dice 0.87) and gave practically the same curve (median $r$ = 0.9999 in A, 0.9998 in B).
The peak was 2\,\% higher and CBV and CBF 2\,\% lower in both planes (Table~\ref{tab:sq1_main}); these
differences are significant because they are very consistent, but they are an order of magnitude smaller than
the Bland--Altman limits of the grown mask and did not differ between planes (all interaction $p_\mathrm{Holm}
\geq 0.09$).

\textbf{Grown full ICA.} The grown mask is about ten times larger (median Dice with \texttt{ica\_top} 0.19)
and samples the ICA upstream. Its AIF arrived 0.15\,s (A) and 0.10\,s (B) earlier and was temporally wider
(temporal SD $+0.22$\,s and $+0.15$\,s), as expected for an average over a longer, more proximal segment
(median $r$ with \texttt{ica\_top} 0.94 and 0.95). The amplitude, however, \emph{changed direction} between
the views: the peak was 7\,\% higher than \texttt{ica\_top} in plane A but 11\,\% lower in plane B
(Fig.~\ref{fig:sq1_forest}; plane interaction $p_\mathrm{Holm} = 6\times10^{-10}$), and the first-pass AUC
was 19\,\% lower in B. The perfusion parameters followed the AIF: in plane B CBV, CBF and MTT were 32\,\%,
22\,\% and 15\,\% higher with the grown AIF (all $p_\mathrm{Holm} < 10^{-6}$), in plane A they did not differ on
average. Even without a mean difference, agreement in plane A was poor for individual runs: the 95\,\% limits of
agreement of CBV were $-67$\,\% to $+181$\,\% (Table~\ref{tab:sq1_main}). Tmax did not differ in either plane
(medians identical at the 0.25\,s frame interval).

\textbf{Why the direction flips.} Of the candidate factors, only vessel overlap was consistently related to the
difference between the two AIFs. The ring around the \texttt{ica\_top} mask contained far more vessel pixels in
the lateral view than in the frontal view (median 0.62 vs.\ 0.34), whereas around the grown mask, which lies
largely on the cervical ICA, it was similar in both views (0.26 vs.\ 0.24; Table~\ref{tab:sq1_context}). The
difference in overlap between the masks correlated with the difference in peak height ($\rho = 0.67$) and in CBV
($\rho = -0.65$, both planes pooled, $n = 99$, $p < 10^{-4}$), and within each plane (peak: $\rho = 0.48$ in A,
$0.36$ in B; CBV: $\rho = -0.61$ and $-0.38$; all $p < 0.01$). Neither the pixel-count ratio nor the size or
spread of the supraclinoid mask (foreshortening) showed such a relation for peak height or CBV. Normalised to
the whole-image contrast, the \texttt{ica\_top} AIF was equal in both views (B/A 1.00, IQR 0.89--1.12,
$p_\mathrm{Holm} = 1.0$), as was \texttt{ica\_full\_model} (0.99), whereas the grown AIF was 16\,\% lower in the
lateral view (0.84, IQR 0.75--0.93, $p_\mathrm{Holm} = 6\times10^{-8}$).

\begin{table*}[t]
\centering
\caption{Effect of the sampling segment: full ICA minus ICA top, per plane, for pairs in which both AIFs are
trusted. Effect = mean (paired $t$) or median ($^\dagger$Wilcoxon) of the difference; amplitudes and perfusion
parameters as \%, times in s. $p$: Holm-corrected within comparison and plane. Bottom row: Bland--Altman 95\,\%
limits of agreement of CBV. All 15 features: Table~\ref{tab:sq1} (Appendix~D).}
\label{tab:sq1_main}
\footnotesize
\begin{tabular}{l rr rr rr rr}
\hline
 & \multicolumn{4}{c}{\texttt{ica\_full\_model} $-$ \texttt{ica\_top}} & \multicolumn{4}{c}{\texttt{ica\_full\_grown} $-$ \texttt{ica\_top}} \\
 & \multicolumn{2}{c}{A ($n$=44)} & \multicolumn{2}{c}{B ($n$=51)} & \multicolumn{2}{c}{A ($n$=46)} & \multicolumn{2}{c}{B ($n$=53)} \\
Feature & effect & $p$ & effect & $p$ & effect & $p$ & effect & $p$ \\
\hline
Arrival time (s)   & $-0.001^\dagger$ & 0.46  & $-0.002^\dagger$ & $<$0.001 & $-0.154^\dagger$ & $<$0.001 & $-0.101^\dagger$ & $<$0.001 \\
Peak height        & $+2.0\%^\dagger$ & $<$0.001 & $+2.1\%^\dagger$ & $<$0.001 & $+7.2\%$  & 0.002 & $-11.1\%$ & $<$0.001 \\
First-pass AUC     & $+1.4\%^\dagger$ & 0.014 & $+1.3\%^\dagger$ & 0.008 & $-0.4\%^\dagger$ & 1.00 & $-18.8\%^\dagger$ & $<$0.001 \\
FWHM (s)           & $-0.001^\dagger$ & 1.00  & $-0.003^\dagger$ & 0.10  & $-0.088$ & 0.007 & $-0.162$ & $<$0.001 \\
Temporal SD (s)    & $-0.001^\dagger$ & 0.74  & $0.000^\dagger$  & 1.00  & $+0.216$ & $<$0.001 & $+0.149$ & $<$0.001 \\
CBV                & $-2.1\%^\dagger$ & 0.007 & $-2.2\%^\dagger$ & $<$0.001 & $-2.1\%^\dagger$ & 1.00 & $+31.5\%^\dagger$ & $<$0.001 \\
CBF                & $-1.8\%^\dagger$ & 0.001 & $-1.8\%^\dagger$ & 0.008 & $+0.6\%^\dagger$ & 1.00 & $+22.1\%$ & $<$0.001 \\
MTT                & $-0.8\%^\dagger$ & 0.052 & $+0.2\%^\dagger$ & 1.00  & $+4.0\%^\dagger$ & 1.00 & $+15.0\%$ & $<$0.001 \\
\hline
CBV 95\,\% LoA     & \multicolumn{2}{c}{$-16$ to $+11\%$} & \multicolumn{2}{c}{$-9$ to $+6\%$} & \multicolumn{2}{c}{$-67$ to $+181\%$} & \multicolumn{2}{c}{$-17$ to $+135\%$} \\
\hline
\end{tabular}
\end{table*}

\begin{table}[t]
\centering
\caption{Surroundings of the masks per view (medians over the grown-vs-top pairs) and the AIF normalised to the
whole-image contrast, plane B/A (all 68 A/B pairs; Section~\ref{sec:plane_methods}).}
\label{tab:sq1_context}
\footnotesize
\begin{tabular}{lccc}
\hline
 & \texttt{ica\_top} & \texttt{ica\_full\_model} & \texttt{ica\_full\_grown} \\
\hline
Mask size A / B (px)                 & 1138 / 1553 & 1181 / 1584$^\ast$ & 12141 / 14380 \\
Vessel fraction in ring, A / B        & 0.34 / 0.62 & 0.34 / 0.57$^\ast$ & 0.24 / 0.26 \\
AIF / whole-image contrast, B/A       & 1.00 & 0.99 & 0.84 \\
\quad IQR; $p_\mathrm{Holm}$          & 0.89--1.12; 1.0 & 0.86--1.13; 1.0 & 0.75--0.93; $6\times10^{-8}$ \\
\hline
\multicolumn{4}{l}{$^\ast$ over the model-vs-top pairs.}
\end{tabular}
\end{table}

\begin{figure*}[t]
\centering
\includegraphics[width=\textwidth]{figures/sq1_forest.png}
\caption{Full ICA minus ICA top for every feature, per plane (blue: A, frontal; orange: B, lateral). Median
difference divided by the SD of the differences, with IQR, so that all features share one axis; ringed:
$p_\mathrm{Holm} < 0.05$. Left: model-segmented full ICA (nearly the same mask as ICA top). Right: grown full ICA,
whose peak, first-pass AUC and perfusion parameters change direction or size between the views.}
\label{fig:sq1_forest}
\end{figure*}

% ---------------------------------------------------------------------
% DISCUSSION / RECOMMENDATION
% ---------------------------------------------------------------------
\subsection{Does the sampling segment matter?}
\label{sec:sq1_discussion}

\textbf{Yes, but only when the segment really changes.} A network trained to segment the full ICA returned
practically the supraclinoid segment again (Dice 0.87, $r \approx 1$), and its AIF differed by about 2\,\%, which
is negligible next to the other sources of variation studied here. Growing the mask along the ICA, by contrast,
changed the AIF in three ways. Two follow directly from sampling further upstream and averaging a longer
segment: earlier arrival and a wider bolus, the delay and dispersion effects known from MR and CT
perfusion~\cite{calamante2000,calamante2013}. Here both are small, the arrival shift (0.10--0.15\,s) is shorter
than one frame interval and Tmax was unaffected. The third is not: the amplitude of the grown AIF relative to
the supraclinoid AIF depends on the view, so that the grown AIF raises CBV, CBF and MTT by 15--32\,\% in the
lateral view and not in the frontal view, with wide limits of agreement in both. Unlike in CT perfusion, where
AIF placement mainly affected CBF and MTT and CBV least~\cite{ferreira2010}, CBV was the most affected
parameter here, which fits an amplitude effect: CBV scales directly with the area under the AIF.

\textbf{Why.} In a projection image, the AIF is the iodine integrated along the X-ray path, so anything that lies
in the path in one view and not in the other changes it. The difference between the segments tracked the vessel
overlap around the supraclinoid mask, which in the lateral view lies amid the siphon loops and their branches,
while pixel count and foreshortening of the supraclinoid segment did not explain it. Normalised to all contrast
in the image, the supraclinoid AIF was the same in both views, and so was the model-segmented AIF; only the
grown AIF was lower in the lateral view. The view dependence therefore lies in the long-segment AIF, most of
whose pixels lie on the petrous and cervical ICA; a plausible but untested explanation is that the projected
path length through these curved segments differs between the frontal and lateral view. This agrees with the rationale of perfDSA for the supraclinoid
segment, which avoids overprojected contrast and contrast stasis in the lower ICA~\cite{su2025perfdsa}. That the
correlation with overlap is clear but moderate ($|\rho| \approx 0.4$--$0.7$) means that overlap is the main
measured factor, not the only one; the analysis has no ground-truth AIF, so it shows which AIF is consistent
across views, not which is correct.

\textbf{Recommendation.}
(i) Sample the AIF from the supraclinoid ICA (\texttt{ica\_top}), as perfDSA does: it is the AIF that is
consistent between views after removing the global contrast scaling, it is closest to the tissue (least delay
and dispersion~\cite{calamante2000}), and a longer segment adds a view-dependent bias of up to a third in CBV.
(ii) A full-ICA segmentation model brings no benefit, since it reproduces the supraclinoid AIF.
(iii) Because the supraclinoid segment is crowded in the lateral view, report the overlap around the AIF mask
(ring vessel fraction, venous fraction) as a quality measure and, when one plane is chosen, prefer the view with
less overlap, which the plane selection of this study does in part by choosing the plane with the lower venous
fraction.
(iv) Keep the sampling segment fixed within a study; perfusion values obtained with different segments, or with
a long segment in different views, should not be compared directly.

\textbf{Limitations.} Pairs required both AIFs to be trusted, so the results describe the effect of the segment
when both extractions are valid; the QC used the tail threshold of 0.30, and the stricter 0.15 supported by the
truncation analysis may change the pair counts slightly. The overlap and foreshortening measures are 2D proxies;
the vessel geometry in 3D was not available.
